%% burst_sabre.tex  —  dSABRE / dSABRE-BurstExt paper
%% Target: DAC / ICCAD 2026  (6 pages) or IEEE TCAD (journal extension)
%% Compile:
%%   pdflatex burst_sabre
%%   bibtex   burst_sabre
%%   pdflatex burst_sabre
%%   pdflatex burst_sabre

\documentclass[conference]{IEEEtran}
\IEEEoverridecommandlockouts

\usepackage{cite}
\usepackage{amsmath,amssymb,amsfonts}
\usepackage{algorithmic}
\usepackage{algorithm}
\usepackage{graphicx}
\usepackage{textcomp}
\usepackage{xcolor}
\usepackage{booktabs}
\usepackage{multirow}
\usepackage{url}
\usepackage{tikz}
\usepackage{pgfplots}
\pgfplotsset{compat=1.18}
\usetikzlibrary{shapes.geometric,arrows,positioning,fit,backgrounds,patterns}

\def\BibTeX{{\rm B\kern-.05em{\sc i\kern-.025em b}\kern-.08em
    T\kern-.1667em\lower.7ex\hbox{E}\kern-.125emX}}

%% ── Notation macros ──────────────────────────────────────────────────────────
\newcommand{\dSABRE}{\textsc{dSABRE}}
\newcommand{\BurstExt}{\textsc{dSABRE-BurstExt}}
\newcommand{\TeleSABRE}{\textsc{TeleSABRE}}
\newcommand{\SABRE}{\textsc{SABRE}}
\newcommand{\eig}[1]{\mathcal{E}(#1)}
\newcommand{\arch}{\mathcal{A}}
\newcommand{\dphys}{d_{\mathrm{phys}}}
\newcommand{\dintra}{d_{\mathrm{intra}}}
\newcommand{\dcore}{d_{\mathrm{core}}}

\begin{document}

\title{dSABRE: A Simple, Effective Router\\
       for Distributed Quantum Circuits}

\author{
  \IEEEauthorblockN{[Author names blinded for review]}
  \IEEEauthorblockA{\textit{[Affiliation blinded for review]}}
}

\maketitle

%%─────────────────────────────────────────────────────────────────────────────
\begin{abstract}
Routing a quantum circuit on a distributed quantum computer (DQC) consumes
one EPR pair for every two-qubit gate whose operands lie in different cores.
We present \dSABRE{}, a \SABRE{}-style distributed router with a five-term
teleportation score and a proactive congestion-relief mechanism, and
\BurstExt{}, a $\sim$50-line subclass that replaces the topological-order
extended set with a \emph{wire-balanced round-robin} lookahead --- adding no
new scoring term, hyperparameter, or auxiliary data structure.
Both routers operate on any DQC modelled as a weighted graph
$G_\arch=(V,E)$ over physical qubits, with cheap intra-core edges and
costly inter-core links; we instantiate the implementation on grid-of-grid
topologies (B-grid, H-grid) for simplicity but the algorithm is
architecture-agnostic.
Across twelve MQT-Bench circuits (25- and 36-qubit) on a $2\!\times\!2$
B-grid, \BurstExt{} reduces EPR consumption over \TeleSABRE{}
\cite{russo2025telesabre} by a geometric mean of \textbf{33.6\%} under
identical initial layouts (\dSABRE{}: 25.3\%); on 64-qubit circuits over a
$2\!\times\!3$ H-grid \BurstExt{} reaches \textbf{28.9\%}, and completes a
Random instance on which \TeleSABRE{} fails to converge.
On the combined routing cost
$3\!\cdot\!\#\mathrm{SWAP}+10\!\cdot\!\#\mathrm{EPR}$, the same routes save
\textbf{17.6\%} on the 25q$+$36q suite and stay competitive at 64q.
EPR cost is counted at one Bell pair per teledata move and one per telegate
(cat-entangler protocol~\cite{eisert2000optimal,yimsiriwattana2004generalized}).
\end{abstract}

\begin{IEEEkeywords}
distributed quantum computing, circuit routing, qubit teleportation,
commutativity, SABRE, EPR minimisation, qubit mapping
\end{IEEEkeywords}

%%─────────────────────────────────────────────────────────────────────────────
\section{Introduction}
\label{sec:intro}

The near-term path to fault-tolerant quantum computing requires processors with
hundreds to thousands of high-fidelity qubits.
No single monolithic chip can economically deliver this scale today.
Instead, the field is converging on \emph{distributed quantum computers}
(DQC)~\cite{monroe2014large,wehner2018quantum,stephenson2020high}: multiple
small quantum processing units (cores) interconnected by photonic or microwave
links that support remote entanglement generation.

Executing a quantum circuit on a DQC requires a \emph{qubit mapping and routing}
pass.
Intra-core two-qubit gates between non-adjacent qubits are resolved with SWAP
insertions; inter-core gates require \emph{quantum teleportation}: one
pre-shared EPR pair is consumed to move a logical qubit across an inter-core
link.
EPR generation is slow (microseconds to milliseconds), noisy, and rate-limited
in current hardware, making EPR minimisation the dominant compilation objective
for DQC.

State-of-the-art distributed routers adapt the \SABRE{}
algorithm~\cite{li2019tackling} to DQC by scoring teleportation candidates using
heuristics that measure how much closer the teleported qubit gets to its
immediate partner gate, with a lookahead term over a fixed window of upcoming
gates~\cite{andres2019automated,ferrari2021compiler,wu2022collcomm,
rodrigues2022noise,baker2020time}.
A fundamental limitation of all these methods is that they evaluate each
teleport in isolation:\ they do not consider whether multiple pending gates on
the same wire would benefit from the same qubit displacement.
In circuits with substantial commutativity---where long runs of gates on a qubit
wire share identical eigen-directions and are mutually commuting---a single
teleport can simultaneously reduce the distance to \emph{many} pending gates.
Ignoring this structure leads to under-valuation of such teleports and, in
practice, to excessive EPR consumption.

\textbf{Contributions.}
We make the following contributions:

\begin{enumerate}
  \item \textbf{\dSABRE{} router.}
    A distributed routing algorithm with a five-term teleportation scoring
    function (preparation cost, capacity penalty, hop gain, front-layer
    distance gain, decay-weighted extended-set gain) and a proactive
    \emph{congestion relief} mechanism that preemptively relocates idle
    qubits out of high-demand cores before they become routing bottlenecks.
    Under identical initial layouts \dSABRE{} achieves a geometric-mean
    \textbf{25.3\%} EPR reduction over \TeleSABRE{}~\cite{russo2025telesabre}
    across 12 benchmark circuits, and completes routing on a 64-qubit
    instance where \TeleSABRE{} fails to converge.

  \item \textbf{Wire-balanced extended set.}
    We identify a structural limitation of \SABRE{}-style lookahead in the
    distributed setting: building the extended set in DAG topological order
    interleaves wires arbitrarily, pushing pending gates that share a wire
    with the current front beyond the lookahead window.
    Replacing topological-order traversal with a round-robin sweep across
    front-layer wires lets the lookahead see the gates that benefit from
    the same teleport, and recovers most of the EPR reduction one would
    obtain from a full commutativity-aware redesign.
    The resulting router, \BurstExt{}, is a $\sim$50-line subclass of
    \dSABRE{} that overrides only the extended-set construction; no new
    scoring term, no new hyperparameter, no auxiliary data structure.
    Under identical initial layouts \BurstExt{} achieves a geometric-mean
    \textbf{33.6\%} EPR reduction over \TeleSABRE{}.

  \item \textbf{Empirical evaluation.}
    Experiments on twelve MQT-Bench circuits~\cite{quetschlich2023mqtbench}
    (six 25-qubit, six 36-qubit) mapped to a $2\!\times\!2$ B-grid of
    $4\!\times\!4$-qubit cores show that \BurstExt{} matches or beats
    \dSABRE{} on 11 of 12 circuits and outperforms \TeleSABRE{} on all
    twelve under identical initial layouts, with no per-circuit tuning.
    A 64-qubit study on a $2\!\times\!3$ H-grid further demonstrates that
    \dSABRE{} completes routing on a Random circuit where \TeleSABRE{}
    fails to converge.
\end{enumerate}

%%─────────────────────────────────────────────────────────────────────────────
\section{Background}
\label{sec:background}

\subsection{Distributed Quantum Architecture}
\label{sec:arch}

We model a DQC architecture $\arch$ abstractly as a weighted graph
$G_\arch=(V,E,w)$ whose vertices are physical qubits.
The vertex set is partitioned into $K$ \emph{cores} $V = V_1\sqcup\cdots\sqcup V_K$;
intra-core edges $E_{\mathrm{intra}}$ describe in-core connectivity and carry
weight~$1$, and inter-core edges $E_{\mathrm{inter}}$ connect designated
\emph{communication ports} on neighbouring cores and carry a much larger
weight $w_{\mathrm{inter}} \gg 1$ reflecting the cost of an EPR pair.
\dSABRE{} treats $G_\arch$ as a black box and only queries it through three
distance tables: physical $\dphys$, intra-core $\dintra$, and core-graph
$\dcore$.
The router therefore extends to \emph{any} DQC topology --- arbitrary
intra-core graphs, mixed core sizes, sparse inter-core meshes, hierarchical
trees of clusters, etc. --- by simply supplying these tables.

For experimental simplicity our reference implementation instantiates
$G_\arch$ as a \emph{grid-of-grids}: each core is an $m\!\times\!m$ 2D
nearest-neighbour grid, and the cores themselves form an $r\!\times\!s$
super-grid joined by a small set of inter-core links between specific
boundary qubits.
We use three port arrangements (Section~\ref{sec:impl}): \emph{basic}
(ports at mid-edges), \emph{B-grid} (staggered ports), and \emph{H-grid}
(asymmetrically skewed ports).

\paragraph{Hierarchical distance precomputation.}
Because intra-core and inter-core costs decouple, the all-pairs distance
$\dphys$ over $G_\arch$ admits a two-level Dijkstra computation that is
substantially cheaper than running Dijkstra on the full $|V|$-vertex graph.
First, all-pairs distances are computed once \emph{inside} a single core
$V_i$ (size $|V_i|=M$) in $O(M^2 \log M)$ time, giving $\dintra$.
Second, the \emph{core graph} (a multigraph of $K$ super-nodes connected
by inter-core links) is solved in $O(K^2 \log K)$ time, giving $\dcore$.
Third, $\dphys(p,q)$ for any $p\in V_i$, $q\in V_j$ is recovered as
$\dintra(p,\pi_i^*) + \dcore(i,j) + \dintra(\pi_j^*,q)$ minimised over
boundary-port pairs $(\pi_i^*, \pi_j^*)$ realising the optimal core-graph
path.
This brings precomputation from $O((KM)^2 \log(KM))$ down to
$O(KM^2\log M + K^2\log K)$ and is the only architecture-dependent setup
\dSABRE{} performs; routing itself uses $O(1)$ table lookups thereafter.

\subsection{Qubit Mapping and Routing}

Given circuit DAG $C$ and initial layout $\ell:\logi{i}\!\mapsto\!\phys{\sigma(i)}$,
routing inserts the minimum number of SWAP and teleportation operations so every
gate can execute on adjacent physical qubits.
Cost model: each local SWAP costs $c_{\mathrm{swap}}$; each inter-core
teleport costs $c_{\mathrm{tele}}$ (one EPR pair).

\subsection{SABRE}

\SABRE{}~\cite{li2019tackling} selects the SWAP minimising:
\begin{equation}
  H = \tfrac{1}{|F|}\!\sum_{g\in F}\!\Delta_F(g)
    + w\,\tfrac{1}{|E|}\!\sum_{g\in E}\!\Delta_E(g),
  \label{eq:sabre}
\end{equation}
where $F$ is the front layer, $E$ the extended lookahead, and $\Delta(g)$ the
gate-distance change induced by the swap.

%%─────────────────────────────────────────────────────────────────────────────
\section{dSABRE}
\label{sec:dsabre}

\dSABRE{} is a distributed routing algorithm we design that extends \SABRE{} to
multi-core architectures.
It replaces SWAP proposals with \emph{teleport proposals}: for each inter-core
gate in the front layer, both endpoint qubits are considered as teleportation
sources, and every reachable communication link on their respective core
boundaries generates a candidate.

\subsection{Teleportation Scoring}

Each candidate is scored as:
\begin{equation}
  s = \underbrace{d_{\mathrm{prep}}}_{\text{staging}} +
      \underbrace{c_{\mathrm{cap}}}_{\text{capacity}} -
      \underbrace{g_{\mathrm{hop}}}_{\text{hop gain}} -
      \underbrace{\Delta_F}_{\text{front gain}} -
      \underbrace{w_e\,\tilde{\Delta}_E}_{\text{lookahead}},
  \label{eq:dsabre}
\end{equation}
where lower scores are preferred.  Each term is defined as follows.

\textbf{Staging cost} $d_{\mathrm{prep}}$ is the number of intra-core SWAPs
needed to move the qubit to the communication port, plus any eviction SWAPs
needed to free occupied comm ports:
\[
  d_{\mathrm{prep}} = d_{\mathrm{intra}}(p_{\mathrm{src}}, p_{\mathrm{ns}})
                    + d_{\mathrm{evict}}(p_{\mathrm{comm,src}})
                    + d_{\mathrm{evict}}(p_{\mathrm{comm,dst}}).
\]

\textbf{Capacity penalty} $c_{\mathrm{cap}} = c_{\mathrm{pen}}\,\max(0,\,\tau - f_{\mathrm{dst}})$
discourages routing into nearly-full cores.
$f_{\mathrm{dst}}$ is the number of free slots in the destination core,
$\tau$ the capacity threshold, and $c_{\mathrm{pen}}$ a cost multiplier.
This prevents progressive crowding of hot cores that leads to deadlock.

\textbf{Hop gain} $g_{\mathrm{hop}} = w_h\,(\dcore(c_{\mathrm{src}},c_{\mathrm{tgt}}) - \dcore(c_{\mathrm{next}},c_{\mathrm{tgt}}))$
rewards moving the qubit closer to the core that holds its gate partner,
acting as a coarse directional bias on the core graph.

\textbf{Front-layer gain} $\Delta_F = \sum_{g \in F, q \in g} (d_{\mathrm{old}}^g - d_{\mathrm{new}}^g)$
sums the distance reduction induced by the teleport over all front-layer
gates involving the moving qubit, mirroring the \SABRE{} objective.

\textbf{Decay-weighted lookahead} $\tilde{\Delta}_E = \sum_{i,g \in E_i} \gamma^i (d_{\mathrm{old}}^g - d_{\mathrm{new}}^g)$
applies an exponential decay factor $\gamma^i$ at lookahead depth $i$, so
that nearby gates in the extended set exert more influence than distant ones.
This softens the sharp boundary between front and extended layers present in
vanilla \SABRE{}.

\subsection{Proactive Congestion Relief}
\label{sec:relief}

In dense layouts, a core can become a bottleneck: many pending gates demand
qubits from it simultaneously, but its communication ports are saturated.
Standard reactive scoring (Eq.~\ref{eq:dsabre}) does not address this until
the bottleneck triggers a deadlock.

\dSABRE{} adds a \emph{proactive congestion relief} mechanism.
At each routing iteration, we compute a demand vector $d[c]$ over all cores
by counting inter-core gates in the lookahead window whose shortest path
traverses core~$c$.
If any core~$c$ satisfies $d[c] \ge \theta_d$ and $f_c \le \theta_f$
(high demand, few free slots), we generate \emph{relief candidates}: teleports
that move the most inactively-used logical qubits \emph{out} of $c$ and into a
less-loaded neighbour, scored by a dedicated relief bonus proportional to the
demand--free-slot imbalance.
These candidates compete directly with gate-driven teleports and are selected
only when they improve the aggregate score.
The mechanism prevents deadlock spirals by redistributing qubit load
\emph{before} the core fills entirely.

\subsection{Deadlock Recovery}

\dSABRE{} maintains a \emph{backup stack}: every $\Delta_{\mathrm{lim}}$
iterations without gate execution, a checkpoint is pushed.
After $N_{\mathrm{bkp}}$ consecutive deadlock events, the router rolls back
to the most recent checkpoint and tries alternative candidates.
This escapes local minima in the scoring landscape without abandoning the
current routing attempt.

\subsection{Comparison with TeleSABRE}

Table~\ref{tab:dsabre_vs_ts} compares \dSABRE{} and \TeleSABRE{} conceptually.

\begin{table}[t]
\centering
\caption{Design comparison: \dSABRE{} vs.\ \TeleSABRE{}}
\label{tab:dsabre_vs_ts}
\setlength{\tabcolsep}{4pt}
\begin{tabular}{@{}lll@{}}
\toprule
Aspect & \dSABRE{} & \TeleSABRE{} \\
\midrule
Language         & Python                 & C++ \\
Operation types  & Teledata (1 EPR)       & Teledata + Telegate (1 EPR each) \\
Scoring          & 5-term heuristic       & Contracted comm-graph energy \\
Lookahead        & Decay-weighted         & Fixed extended set \\
Congestion       & Proactive relief       & Implicit via capacity terms \\
Initial layout   & Locality-aware         & Hungarian + opt.\ init.\ \\
Deadlock escape  & Backup/rollback        & Safety-valve iterations \\
\bottomrule
\end{tabular}
\end{table}

The key architectural difference is that \TeleSABRE{} evaluates candidates
using a \emph{contracted communication-graph energy}: the circuit is projected
onto the core graph and a global energy function measures how far all logical
qubits are from their nearest gate partners.
This provides a broader view but does not account for intra-core staging cost
or core capacity.
\dSABRE{} instead uses a local, gate-centric score that explicitly models
staging, capacity, and directional progress, combined with proactive load
balancing.
Empirically (Section~\ref{sec:telesabre_comp}), \dSABRE{}'s routing heuristic
produces fewer EPRs than \TeleSABRE{} on 4 of 6 circuits under identical
initial layouts, demonstrating that the scoring difference---not layout---drives
the advantage.

%%─────────────────────────────────────────────────────────────────────────────
\section{Wire-Balanced Extended Set}
\label{sec:burstext}

\subsection{The Lookahead Bias}

\SABRE{} (Eq.~\ref{eq:sabre}) and \dSABRE{} (Eq.~\ref{eq:dsabre}) both rely
on an \emph{extended-set} term $E$ that estimates how a candidate action
affects gates beyond the immediate front layer.
The extended set is conventionally built by walking the circuit DAG in
topological order and collecting the next $|E|$ unscheduled 2-qubit gates.

This construction has a subtle but consequential bias in the distributed
setting.
Consider a front-layer gate $g_F$ on qubits $(u,v)$ in different cores, and
suppose wire $u$ continues with a long run of mutually commuting 2-qubit
gates $g_1, g_2, \ldots, g_k$ that all share the same partner core as $g_F$.
A teleport that moves $u$ closer to its partner core would simultaneously
reduce the distance for \emph{all} of $g_F, g_1, \ldots, g_k$.
This is precisely the regime where teleportation amortises well across many
gates --- and where good scoring matters most.

Yet in the topological-order extended set, $g_1, \ldots, g_k$ are typically
\emph{not} present.
Topological order interleaves gates from many wires; long before the walk
returns to wire $u$, it has filled the lookahead window with unrelated gates
on other wires.
The scoring function therefore sees only $g_F$ and a handful of unrelated
front-adjacent gates, and treats the teleport identically to one that would
benefit only $g_F$ in isolation.
The result is systematic under-valuation of teleports that lie at the head
of long commutativity chains.

\subsection{Round-Robin Construction}

Our fix is to replace the topological walk with a wire-balanced round-robin.
Let $Q_F = \bigcup_{g\in F}\,g.\mathit{qargs}$ be the set of qubits appearing
in the front layer.
We bucket every non-front 2-qubit gate by each of its qubits and then sweep
across wires in round-robin order, picking the next pending gate from each
wire on each pass, until the extended set is filled
(Algorithm~\ref{alg:wirebal}).

\begin{algorithm}[t]
\caption{\texttt{ExtendedSet-Wire}$(C, F, |E|)$}
\label{alg:wirebal}
\begin{algorithmic}[1]
\STATE $Q_F \leftarrow \bigcup_{g\in F}\,g.\mathit{qargs}$
\STATE $B[q] \leftarrow []$ for each $q \in Q_F$
\FOR{each non-front 2-qubit gate $g$ in topological order}
  \FOR{each $q \in g.\mathit{qargs} \cap Q_F$}
    \STATE append $g$ to $B[q]$
  \ENDFOR
\ENDFOR
\STATE $E \leftarrow []$;\ $S \leftarrow \emptyset$;\
       cursor$[q] \leftarrow 0$;\ active $\leftarrow Q_F$
\WHILE{active $\ne \emptyset$ \AND $|E| < |E|_{\max}$}
  \STATE next $\leftarrow$ \texttt{[]}
  \FOR{$q \in $ active}
    \STATE advance cursor$[q]$ past gates already in $S$
    \IF{cursor$[q] < |B[q]|$}
      \STATE $g \leftarrow B[q][\,\text{cursor}[q]\,]$;\
             append $g$ to $E$;\ insert $g$ into $S$
      \STATE cursor$[q] \mathrel{+}= 1$;\ append $q$ to next
      \IF{$|E| \ge |E|_{\max}$} \textbf{break}\ \ENDIF
    \ENDIF
  \ENDFOR
  \STATE active $\leftarrow$ next
\ENDWHILE
\RETURN $E$
\end{algorithmic}
\end{algorithm}

The cost is $O(m)$ to bucket the gates plus $O(|E|)$ to sweep, both
proportional to inputs the scoring loop already touches.
No new state is maintained between iterations.

\subsection{Running Example}
\label{sec:running}

Consider a six-qubit circuit on a two-core DQC with
$C_L=\{q_0,q_1,q_2\}$, $C_R=\{q_3,q_4,q_5\}$, and one inter-core link
between $q_2$ and $q_3$.
The circuit is
\begin{center}\small
$g_F{:}\,\mathrm{CX}(q_0,q_3)$\quad\quad
$g_1{:}\,\mathrm{CX}(q_4,q_5)$\quad
$g_2{:}\,\mathrm{CX}(q_1,q_2)$\\[2pt]
$g_3{:}\,\mathrm{CX}(q_4,q_5)$\quad
$g_4{:}\,\mathrm{CX}(q_1,q_2)$\quad
$g_5{:}\,\mathrm{CX}(q_0,q_4)$\\[2pt]
$g_6{:}\,\mathrm{CX}(q_0,q_5)$\quad\quad
$g_7{:}\,\mathrm{CX}(q_0,q_3)$
\end{center}
Here $F=\{g_F\}$ is the only inter-core front-layer gate
($Q_F=\{q_0,q_3\}$), and the gates $g_F, g_5, g_6, g_7$ form a
\emph{commutativity chain} on wire $q_0$:
they all use $q_0$ as control, so they share the eigen-direction
$\{|0\rangle,|1\rangle\}$ on that wire and pairwise commute on
$q_0$~\cite{hong2024commutativity}.
A single teleport moving $q_0$ from $C_L$ to $C_R$ would resolve the
inter-core distance of \emph{every} gate in this chain --- the
canonical regime where teleportation amortises.
Suppose $|E|_{\max}=4$.

\textbf{Topological-order extended set.}
After $g_F$ leaves the front, the next four 2-qubit gates in DAG order
are the wires currently \emph{ready}:
$E_{\mathrm{topo}} = [g_1, g_2, g_3, g_4]$.
None of them touches $q_0$; the entire $q_0$-control chain
$g_5,g_6,g_7$ is pushed beyond the lookahead window by the
unrelated $C_R$-internal and $C_L$-internal traffic on wires $q_4,q_5$
and $q_1,q_2$.
The scoring of ``teleport $q_0$ to $C_R$'' therefore receives
$\tilde{\Delta}_E = 0$ from the lookahead --- the candidate is judged on
$g_F$ alone, exactly as if no commutativity were present.

\textbf{Wire-balanced extended set.}
Bucketing the non-front 2-qubit gates by their front-layer wires yields
$B[q_0] = [g_5, g_6, g_7]$ and $B[q_3] = [g_7]$.
The round-robin sweep emits one gate per active wire per pass:
pass~1 returns $\{g_5, g_7\}$, pass~2 returns $\{g_6\}$ (wire $q_3$
exhausted), and the cursor on $q_0$ skips $g_7$ which is already in $E$.
The result is
$E_{\mathrm{wire}} = [g_5, g_7, g_6]$ --- the \emph{full} $q_0$
commutativity chain.
``Teleport $q_0$ to $C_R$'' now accumulates
$\Delta_F(g_F) + \tilde{\Delta}_E(g_5) + \tilde{\Delta}_E(g_6) +
 \tilde{\Delta}_E(g_7)$, four positive contributions instead of one,
correctly identifying that one EPR pair amortises across four
inter-core gates.
With the topological set, the router would typically commit a separate
teleport for each of $g_5$, $g_6$, $g_7$ as they reach the front layer.

\subsection{Why It Captures Commutativity For Free}

We do not detect commutativity explicitly; we never partition the circuit
into commuting clusters and never test whether any pair of gates shares an
eigen-direction~\cite{hong2024commutativity}.
The wire-balanced extended set captures commutativity \emph{indirectly}:
when a wire has a long run of pending gates with the same partner-core
structure, those gates appear consecutively in $B[q]$ and are pulled into
$E$ on successive round-robin passes.
The distance-gain term $\Delta_E$ in Eq.~\ref{eq:dsabre} then sums their
contributions, automatically up-weighting teleports that benefit many of
them simultaneously, achieved at zero cost in lookahead complexity, scoring
complexity, or hyperparameter tuning.

The resulting router, \BurstExt{}, is implemented as a \texttt{\_extended\_2q}
override on \dSABRE{} totalling roughly fifty lines.
All other components --- proactive congestion relief, decay-weighted
$\Delta_E$, backup/rollback, capacity penalties --- remain unchanged.
Whether explicit commutativity machinery (e.g.\ a hypergraph partitioning
of each wire into maximal commuting runs~\cite{hong2024commutativity},
plus a dedicated scoring term) yields further gains beyond what the
wire-balanced extended set captures implicitly is the subject of a
companion paper.

\subsection{Gate Commutativity (Background)}

For completeness, two gates $g_1, g_2$ on qubit $q$ commute on wire $q$ iff
$\eig{g_1,q} = \eig{g_2,q}$, where $\eig{g,q}$ is the set of simultaneous
eigenstates of the gate's action on qubit $q$~\cite{hong2024commutativity}.
All CNOT controls share direction $\{|0\rangle,|1\rangle\}$, so CNOT gates
sharing a control qubit commute on the control wire; $R_z$ gates all share
$\{|0\rangle,|1\rangle\}$.
The wire-balanced extended set is agnostic to these details: it merely
preserves the option for the scoring function to sum distance gains over
gates that happen to share a wire with the front layer.


%%─────────────────────────────────────────────────────────────────────────────
\section{Experimental Evaluation}
\label{sec:experiments}

\subsection{Setup}

\textbf{Benchmarks.}
We use two circuit sets from the MQT Bench suite~\cite{quetschlich2023mqtbench},
both in the IBM native-gate set (Qiskit opt3 transpilation, measurements and
barriers removed).

\emph{Suite~1 (25-qubit):} quantum amplitude estimation (AE), quantum Fourier
transform (QFT), quantum neural network (QNN), random circuit (Random), GHZ
state preparation, and graph-state preparation.

\emph{Suite~2 (36-qubit):} six additional circuits chosen to span a broader
range of algorithmic families and CX densities: Bernstein--Vazirani (BV, linear
chain), Deutsch--Jozsa (DJ, oracle), W-state preparation (cascaded star),
VQE with SU(2) ansatz (layered variational), QAOA (dense all-to-all), and
quantum phase estimation (QPEexact, hierarchical).

Circuit properties are summarised in Table~\ref{tab:circuits}.

\begin{table}[t]
\centering
\caption{Benchmark circuit properties (measurements/barriers stripped)}
\label{tab:circuits}
\setlength{\tabcolsep}{5pt}
\begin{tabular}{@{}llrrrr@{}}
\toprule
Suite & Circuit    & Qubits & CX gates & Depth & CX/qubit \\
\midrule
\multirow{6}{*}{25q}
 & AE         & 25 &  558 & 395 & 22.3 \\
 & QFT        & 25 &  580 & 173 & 23.2 \\
 & QNN        & 25 & 1223 & 259 & 48.9 \\
 & Random     & 25 & 1124 & 589 & 45.0 \\
 & GHZ        & 25 &   24 &  27 &  1.0 \\
 & Graphstate & 25 &   25 &  19 &  1.0 \\
\midrule
\multirow{6}{*}{36q}
 & BV         & 36 &   17 &  23 &  0.5 \\
 & DJ         & 36 &   35 &  41 &  1.0 \\
 & W-state    & 36 &   70 & 145 &  1.9 \\
 & VQE-SU2   & 36 &  105 &  56 &  2.9 \\
 & QPEexact   & 36 & 1019 & 347 & 28.3 \\
 & QAOA       & 36 & 1200 & 256 & 33.3 \\
\bottomrule
\end{tabular}
\end{table}

\textbf{Architecture.}
All experiments use a $2\!\times\!2$ B-grid architecture: four cores each with
a $4\!\times\!4$ grid of qubits (64 physical slots total), with staggered
communication ports as described in Section~\ref{sec:impl}.

\textbf{Baselines.}
We compare against two baselines: (i)~\dSABRE{}~\cite{andres2019automated}
under identical layout strategies and hardware cost parameters; and
(ii)~\TeleSABRE{}~\cite{russo2025telesabre}, a state-of-the-art distributed
router that jointly optimises intra-core SWAPs, teledata (qubit moves), and
telegate (remote CX execution) operations.
\TeleSABRE{} is run using its compiled C binary with the \texttt{optimize\_initial}
flag enabled (Hungarian assignment + 3-success sampling) on the same
$2\!\times\!2$ B-grid device.

\textbf{Parameters.}
Table~\ref{tab:params} lists all hyperparameters.
Three trials are run per (router, layout) pair; we report the best EPR count
across trials (lower = better), following established practice~\cite{li2019tackling}.
Layout passes: 2 (forward–backward–forward).

\begin{table}[t]
\centering
\caption{Hardware cost and heuristic parameters}
\label{tab:params}
\begin{tabular}{@{}lll@{}}
\toprule
Parameter & Symbol & Value \\
\midrule
SWAP cost             & $c_{\mathrm{swap}}$ & 3 \\
Teleport (EPR) cost   & $c_{\mathrm{tele}}$ & 10 \\
Capacity threshold    & $\tau$              & 3 \\
Capacity penalty      & $c_{\mathrm{pen}}$  & 15 \\
Hop-gain weight       & $w_h$               & 5 \\
Extended-set weight   & $w_e$               & 0.5 \\
Lookahead decay       & $\gamma$            & 0.9 \\
Deadlock limit        & $\Delta_{\rm lim}$  & 50 \\
Max rollbacks         & $N_{\rm bkp}$       & 50 \\
\bottomrule
\end{tabular}
\end{table}

\subsection{Main Results}

Our headline comparison fixes the initial layout (to \TeleSABRE{}'s
\texttt{optimize\_initial} Hungarian assignment) and then runs both
\dSABRE{} and \BurstExt{}, isolating the contribution of the routing
heuristic itself.
Section~\ref{sec:telesabre_comp} reports those results
(Table~\ref{tab:telesabre}, Table~\ref{tab:telesabre_36q}).

\subsection{Comparison with TeleSABRE}
\label{sec:telesabre_comp}

\TeleSABRE{}~\cite{russo2025telesabre} supports three operation types:
intra-core SWAPs, teledata (qubit movement, 1~EPR), and telegate
(remote CX execution in-place via the cat-entangler protocol, 1~EPR).
A naive comparison using each router's own initial layout would conflate
layout quality with routing quality.
To isolate the routing heuristic, we run all routers from
\emph{exactly TeleSABRE's initial layout} (Hungarian assignment,
\texttt{optimize\_initial}=true) as the fixed starting point, followed by
the same two forward/backward passes.
The corresponding configurations are denoted \dSABRE$^{\dagger}$ and
\BurstExt$^{\dagger}$.
We measure each route on three axes:
(i)~EPR pairs (teledata+telegate, 1~Bell pair per primitive);
(ii)~local intra-core SWAPs (\#SWAP);
(iii)~combined routing cost
$\mathrm{Cost}=c_{\mathrm{swap}}\,\#\mathrm{SWAP}+c_{\mathrm{tele}}\,\#\mathrm{EPR}$
with $c_{\mathrm{swap}}=3$ and $c_{\mathrm{tele}}=10$
(Table~\ref{tab:params}).

\begin{table*}[t]
\centering
\caption{Fair comparison on 25-qubit circuits under identical Hungarian initial
  layouts.  EPR$=$teledata$+$telegate (1 Bell pair per primitive);
  Cost $= 3\!\cdot\!\#\mathrm{SWAP}+10\!\cdot\!\#\mathrm{EPR}$.
  $\Delta$ rows: \BurstExt$^{\dagger}$ vs.\ \TeleSABRE{}; geometric mean.
  Best Cost per row in \textbf{bold}.}
\label{tab:telesabre}
\setlength{\tabcolsep}{5pt}
\begin{tabular}{@{}l rrr rrr rrr rr@{}}
\toprule
 & \multicolumn{3}{c}{\TeleSABRE}
 & \multicolumn{3}{c}{\dSABRE$^{\dagger}$}
 & \multicolumn{3}{c}{\BurstExt$^{\dagger}$}
 & \multicolumn{2}{c}{$\Delta$ (\%)} \\
\cmidrule(lr){2-4}\cmidrule(lr){5-7}\cmidrule(lr){8-10}\cmidrule(lr){11-12}
Circuit & EPR & \#SWAP & Cost & EPR & \#SWAP & Cost & EPR & \#SWAP & Cost & EPR & Cost \\
\midrule
AE         &  26 &  292 & \textbf{1136} &  46 &  490 & 1930 &  26 &  390 & 1430 & $\phantom{-}0.0$ & $+25.9$ \\
QFT        &  46 &  405 & \textbf{1675} &  50 &  468 & 1904 &  45 &  435 & 1755 & $-2.2$ & $+4.8$ \\
QNN        &  54 &  659 & \textbf{2517} &  97 &  974 & 3892 &  54 &  902 & 3246 & $\phantom{-}0.0$ & $+29.0$ \\
Random     & 271 & 1248 & 6454 & 203 & 1399 & \textbf{6227} & 199 & 1443 & 6319 & $-26.6$ & $-2.1$ \\
GHZ        &   4 &   36 & 148  &   2 &   15 & \textbf{65}   &   2 &   15 & \textbf{65}  & $-50.0$ & $-56.1$ \\
Graphstate &  13 &   47 & 271  &  10 &   74 & 322  &   8 &   38 & \textbf{194} & $-38.5$ & $-28.4$ \\
\midrule
\textbf{Gmean $\Delta$ vs.\ TS} & & & & ($-0.1$) & & ($+7.0$) & & & & $\mathbf{-22.2}$ & $\mathbf{-10.2}$ \\
\bottomrule
\end{tabular}
\end{table*}

Under identical initial layouts, \BurstExt$^{\dagger}$ outperforms or matches
\TeleSABRE{} on all six 25-qubit circuits, achieving a geometric-mean EPR
reduction of \textbf{22.2\%}.
The gains are largest on GHZ ($-50.0\%$), Graphstate ($-38.5\%$), and
Random ($-26.6\%$);
on AE and QNN \BurstExt{} matches \TeleSABRE{} exactly while improving over
the topological-extended-set baseline (\dSABRE$^{\dagger}$) by 43--44\%.

In combined-cost terms, the picture is more conservative:
\BurstExt$^{\dagger}$ saves \textbf{10.2\%} over \TeleSABRE{} on this 25-qubit
suite (gmean), with strong wins on GHZ ($-56.1\%$) and Graphstate ($-28.4\%$)
balanced by AE ($+25.9\%$), QNN ($+29.0\%$), and QFT ($+4.8\%$):
on the high-CX-density structured circuits, \TeleSABRE{}'s tighter SWAP
budget (e.g.\ 292 vs.\ 390 SWAPs on AE) dominates EPR savings under the
$3{:}10$ weighting.

\textbf{\dSABRE{} vs.\ \TeleSABRE{}.}
The \dSABRE$^{\dagger}$ column isolates \dSABRE{}'s scoring heuristic from
layout quality.
Under TeleSABRE's Hungarian initial layout, \dSABRE{} beats \TeleSABRE{} on
3 of 6 circuits (Random, GHZ, Graphstate), and the suite-level geometric
mean is essentially flat (\textbf{$-0.1\%$})---\dSABRE{}'s scoring breaks
even with \TeleSABRE{} on this small suite, with the headline 25q gains
arriving from the wire-balanced extended set (\BurstExt{}).
The QNN circuit is the one case where \TeleSABRE{} retains an advantage
($97$ vs.\ $54$); its telegate operations (remote CX, 1 EPR via cat-entangler) reduce
per-gate movement for this high-CX-density structured circuit.
On the 64-qubit Random circuit, however, \TeleSABRE{} fails to converge
entirely while \dSABRE{} completes routing with 834~EPRs, demonstrating that
\dSABRE{}'s backup/rollback deadlock-recovery offers stronger completion
guarantees than \TeleSABRE{}'s iteration-capped solver on irregular
interaction patterns.

The \dSABRE$^{\dagger}$ vs.\ \BurstExt$^{\dagger}$ comparison is the cleanest
ablation in the paper: the only code difference between the two routers is
the extended-set construction (Algorithm~\ref{alg:wirebal}).
On QNN this single change cuts EPRs from 97 to 54 ($-44.3\%$); on AE from
46 to 26 ($-43.5\%$); on Random from 203 to 199 ($-2.0\%$).
The wire-balanced lookahead is the dominant lever.

\textbf{36-qubit results.}
Table~\ref{tab:telesabre_36q} extends the comparison to six 36-qubit
circuits on the same B-grid device, chosen to span a broader range of
algorithmic families and CX densities than the original suite.
\BurstExt$^{\dagger}$ outperforms \TeleSABRE{} on all six circuits,
achieving a geometric-mean \textbf{43.3\%} EPR reduction and a combined-cost
reduction of \textbf{24.4\%}.
The EPR gains range from $-18.7\%$ (QAOA, dense all-to-all) to $-53.8\%$
(W-state); cost gains range from $+7.7\%$ (QAOA, where extra SWAPs offset
EPR savings) to $-42.8\%$ (W-state).

Across both suites (12 circuits total), the combined geometric-mean EPR
reduction of \BurstExt$^{\dagger}$ over \TeleSABRE{} is \textbf{33.6\%}
(\dSABRE$^{\dagger}$: \textbf{25.3\%}); the corresponding combined-cost
reductions are \textbf{17.6\%} and \textbf{7.8\%}.
The wire-balanced extended set is therefore responsible for roughly
8~percentage points of additional EPR reduction (and 10~points of cost
reduction) beyond \dSABRE{} on top of the baseline scoring heuristic --- at
the price of one method override and zero extra hyperparameters.

\begin{table*}[t]
\centering
\caption{Fair comparison on 36-qubit circuits (B-grid $2\!\times\!2$ $4\!\times\!4$).
  Same methodology as Table~\ref{tab:telesabre}.
  Best Cost per row in \textbf{bold}.}
\label{tab:telesabre_36q}
\setlength{\tabcolsep}{5pt}
\begin{tabular}{@{}l rrr rrr rrr rr@{}}
\toprule
 & \multicolumn{3}{c}{\TeleSABRE}
 & \multicolumn{3}{c}{\dSABRE$^{\dagger}$}
 & \multicolumn{3}{c}{\BurstExt$^{\dagger}$}
 & \multicolumn{2}{c}{$\Delta$ (\%)} \\
\cmidrule(lr){2-4}\cmidrule(lr){5-7}\cmidrule(lr){8-10}\cmidrule(lr){11-12}
Circuit & EPR & \#SWAP & Cost & EPR & \#SWAP & Cost & EPR & \#SWAP & Cost & EPR & Cost \\
\midrule
BV       &   6 &   39 & 177  &   3 &   27 & \textbf{111} &   3 &   27 & \textbf{111} & $-50.0$ & $-37.3$ \\
DJ       &   4 &   63 & 229  &   3 &   50 & \textbf{180} &   3 &   50 & \textbf{180} & $-25.0$ & $-21.4$ \\
W-state  &  13 &   79 & 367  &   6 &   65 & 255  &   6 &   50 & \textbf{210} & $-53.8$ & $-42.8$ \\
VQE-SU2  &  28 &  134 & 682  &  11 &  156 & 578  &  13 &  137 & \textbf{541} & $-53.6$ & $-20.7$ \\
QPEexact & 190 &  883 & 4549 & 101 &  878 & 3644 &  97 &  849 & \textbf{3517} & $-48.9$ & $-22.7$ \\
QAOA     & 268 & 1165 & \textbf{6175} & 225 & 1474 & 6672 & 218 & 1490 & 6650 & $-18.7$ & $+7.7$ \\
\midrule
\textbf{Gmean $\Delta$ vs.\ TS} & & & & ($-44.1$) & & ($-20.6$) & & & & $\mathbf{-43.3}$ & $\mathbf{-24.4}$ \\
\bottomrule
\end{tabular}
\end{table*}

\textbf{Where the wire-balanced extended set helps and where it does not.}
\BurstExt{} matches or beats \dSABRE{} on 11 of 12 circuits.
The single regression is VQE-SU2 (13 vs.\ 11): VQE's layered-rotation
structure happens to admit a topological-order lookahead that already
captures the relevant inter-layer gates, so the round-robin sweep introduces
slight noise.
The largest reductions over \dSABRE{} appear on AE ($-43.5\%$), Graphstate
($-20.0\%$), QNN ($-44.3\%$), QFT ($-10.0\%$), QAOA ($-3.1\%$), and
QPEexact ($-4.0\%$): circuits where front-layer wires continue with long
runs of pending 2-qubit gates that the topological-order extended set
displaces beyond the lookahead window.

\subsection{Compile Time}
\label{sec:timing}

Table~\ref{tab:timing} reports wall-clock compilation time (single routing
pass, fixed Hungarian initial layout, seed~0) for all four routers on the
25-qubit benchmark circuits.
\TeleSABRE{} is a compiled C++ binary; the dSABRE family is pure-Python.

\begin{table}[t]
\centering
\caption{Compilation time in seconds (25-qubit suite).
  Slowdowns are relative to \TeleSABRE{}.
  \BurstExt\ adds essentially no overhead over \dSABRE.}
\label{tab:timing}
\setlength{\tabcolsep}{4pt}
\begin{tabular}{@{}lrrrr@{}}
\toprule
Circuit & CX & TS (s) & dS (s) & dSE (s) \\
\midrule
AE         &  558 & 0.05 &  0.84 &  0.86 \\
GHZ        &   24 & 0.01 &  0.01 &  0.01 \\
Graphstate &   25 & 0.02 &  0.01 &  0.01 \\
QFT        &  580 & 0.07 &  0.75 &  0.78 \\
QNN        & 1223 & 0.09 &  2.92 &  2.95 \\
Random     & 1124 & 0.38 &  5.24 &  5.41 \\
\bottomrule
\end{tabular}
\end{table}

\TeleSABRE{}'s C++ implementation finishes in under $0.4$\,s on all circuits.
\dSABRE{} is $1$--$34\times$ slower than \TeleSABRE{} ($0.01$--$5.2$\,s), a
gap attributable to Python interpreter overhead rather than algorithmic
complexity.
\BurstExt{} adds essentially no overhead over \dSABRE{}: its only difference
is the extended-set construction, whose cost is $O(m+|E|)$ --- the same order
as the topological-order construction it replaces.

%%─────────────────────────────────────────────────────────────────────────────
\section{Implementation Notes}
\label{sec:impl}

\dSABRE{} and \BurstExt{} are implemented in Python using
NetworkX~\cite{hagberg2008exploring} and Qiskit~\cite{qiskit2024};
\BurstExt{} is a $\sim$50-line subclass overriding only the extended-set
construction method.
The reference implementation uses the three grid-of-grid topologies
introduced in Section~\ref{sec:arch}: \emph{basic} (ports at mid-edges),
\emph{B-grid} (staggered ports), and \emph{H-grid} (skewed ports).
Hierarchical Dijkstra precomputation (Section~\ref{sec:arch}) takes
$\sim$50\,ms for a $2\!\times\!2$ grid of $4\!\times\!4$ cores
($64$~nodes); for a $4\!\times\!4$ grid of $8\!\times\!8$ cores
($1024$ nodes) the $\dphys$ table grows to ${\approx}10^6$ entries
(${\sim}56$\,MB), still tractable.
Beyond ${\sim}1600$ nodes, lazy/sparse alternatives to the precomputed
table should be considered.

%%─────────────────────────────────────────────────────────────────────────────
\subsection{Scalability and 64-Qubit Results}
\label{sec:scalability}

We evaluated the routers at the next architecture size --- a $2\!\times\!3$
H-grid of $4\!\times\!4$-qubit cores (96 physical slots) running 64-qubit
MQT circuits.
One implementation issue had to be addressed first: when all comm-port
slots in a core are simultaneously occupied, \texttt{\_evict} finds no
free intra-core slot; we now skip the teleport gracefully rather than
asserting.
A stronger remedy that routes an intra-core SWAP to free the port is left
for future work.

\begin{table*}[t]
\centering
\caption{64-qubit benchmark on H-grid $2\!\times\!3$ $4\!\times\!4$
under \TeleSABRE{}'s Hungarian initial layout.
``---'' denotes routing failure (TS: solver did not converge within
its iteration budget; \dSABRE-family: routing aborted within 50k iterations).
$\Delta$ rows: \BurstExt$^{\dagger}$ vs.\ \TeleSABRE; geometric mean over the
five circuits TS converged on. Best Cost per row in \textbf{bold}.}
\label{tab:64q}
\setlength{\tabcolsep}{4pt}
\begin{tabular}{@{}l r rrr rrr rrr rr@{}}
\toprule
 & & \multicolumn{3}{c}{\TeleSABRE}
 & \multicolumn{3}{c}{\dSABRE$^{\dagger}$}
 & \multicolumn{3}{c}{\BurstExt$^{\dagger}$}
 & \multicolumn{2}{c}{$\Delta$ (\%)} \\
\cmidrule(lr){3-5}\cmidrule(lr){6-8}\cmidrule(lr){9-11}\cmidrule(lr){12-13}
Circuit & CX & EPR & \#SWAP & Cost & EPR & \#SWAP & Cost & EPR & \#SWAP & Cost & EPR & Cost \\
\midrule
AE         & 1962 & 353 & 1494 & \textbf{8012} & 273 & 2185 & 9285 & 202 & 1717 & \textbf{7171} & $-42.8$ & $-10.5$ \\
GHZ        &   63 &  16 &  132 &  556  &  10 &  130 & \textbf{490}  &  10 &  130 & \textbf{490}  & $-37.5$ & $-11.9$ \\
Graphstate &   64 &  62 &  281 & 1463 &  54 &  292 & \textbf{1416} &  57 &  311 & 1503 & $-8.1$  & $+2.7$  \\
QFT        & 1966 & 312 & 1711 & \textbf{8253} & 301 & 2165 & 9505 & 250 & 2087 & 8761 & $-19.9$ & $+6.2$  \\
QNN        & 8126 & 889 & 5471 & \textbf{25303}& 990 & 8520 & 35460& 613 & 7720 & 29290& $-31.0$ & $+15.8$ \\
Random$^\ast$ & 1627 & --- & --- & --- & \textbf{834} & 4397 & \textbf{21531}& 937 & 4534 & 22972 & --- & --- \\
\midrule
\textbf{Gmean $\Delta$ vs.\ TS} & & & & & ($-14.7$) & & ($+9.8$) & & & & $\mathbf{-28.9}$ & $\mathbf{-0.1}$ \\
\bottomrule
\end{tabular}
\smallskip\\
\footnotesize $^{\ast}$\TeleSABRE{} produces no layout for the 64q Random
circuit; both \dSABRE$^{\dagger}$ and \BurstExt$^{\dagger}$ are therefore run
from the \emph{same} locality-aware initial layout (3 seeds; all three
converged to identical metrics for both routers, so the row is excluded from
the geometric mean).
\end{table*}

The most striking 64-qubit finding is on the Random circuit:
\textbf{\TeleSABRE{} fails to produce a valid routing} (its energy-based
solver does not converge within the configured iteration budget), whereas
\dSABRE{} successfully completes routing with $834$~EPRs (4397~SWAPs,
cost~$21{,}531$) and \BurstExt{} with $937$~EPRs (4534~SWAPs, cost~$22{,}972$).
Because \TeleSABRE{} produces no layout for this circuit, both
\dSABRE{} and \BurstExt{} are run from a \emph{shared} locality-aware
initial layout (identical seed-to-seed, so the 834~vs.\ 937 gap is purely
a routing-heuristic effect, not a layout artefact).
This demonstrates that \dSABRE{}'s backup/rollback deadlock-recovery offers
stronger completion guarantees than \TeleSABRE{}'s safety-valve approach on
circuits with irregular, non-structured interaction patterns.
The wire-balanced extended set actually \emph{regresses} on Random under the
matched initial layout ($937$ vs.\ $834$), consistent with the fact that
uniformly-random gate sequences contain almost no commutativity structure
for the round-robin sweep to exploit; the topological-order extended set
serves \dSABRE{} better in this regime.

On the five circuits where the same-layout comparison is well-defined,
\BurstExt$^{\dagger}$ delivers a geometric-mean \textbf{28.9\%} EPR
reduction over \TeleSABRE{} (and \dSABRE$^{\dagger}$ delivers \textbf{14.7\%}).
Compared to \dSABRE$^{\dagger}$, the wire-balanced extended set yields the
largest EPR gains on the high-CX structured circuits (AE: $-26.0\%$,
QFT: $-16.9\%$, QNN: $-38.1\%$) and ties on GHZ; the regressions are
Graphstate ($+5.6\%$, $54 \to 57$) and Random ($+12.4\%$, $834 \to 937$),
both circuits with low or no commutativity structure.
On combined cost the picture flattens: \BurstExt$^{\dagger}$ is essentially
flat with \TeleSABRE{} (geometric-mean $-0.1\%$ over the same five circuits),
saving $10$--$12\%$ on AE and GHZ but paying $+6$--$+16\%$ on QFT and QNN
where the higher SWAP count outweighs the EPR savings under the $3{:}10$
weighting.  \dSABRE{} is more expensive than \TeleSABRE{} on cost
($+9.8\%$) at this scale.
This sharpens the picture: \BurstExt{} is the better default for
EPR-dominated regimes --- structured workloads where remote-entanglement
generation is the bottleneck --- while \TeleSABRE{} retains an edge when
intra-core SWAPs and EPRs are weighted comparably and CX density is high;
\dSABRE{} remains the more reliable choice for unstructured circuits and
as a robustness fallback.

%%─────────────────────────────────────────────────────────────────────────────
\subsection{Cost-Model Sensitivity}
\label{sec:sensitivity}

The combined-cost ratio $c_{\mathrm{tele}}{:}c_{\mathrm{swap}}$ is
hardware-dependent: near-term microwave-coupled modular superconducting
qubits target ${\sim}10{:}3$~\cite{cuomo2020towards}, while photonic-link
architectures have measured EPR-generation times three to four orders of
magnitude slower than local gates~\cite{stephenson2020high,daiss2021quantum},
pushing the ratio toward $100{:}1$ or beyond.
Table~\ref{tab:sensitivity} sweeps $c_{\mathrm{tele}}$ from~3 to~100
with $c_{\mathrm{swap}}=3$ held fixed, reporting geometric-mean cost
reduction of \dSABRE$^{\dagger}$ and \BurstExt$^{\dagger}$ vs.\
\TeleSABRE{}.

\begin{table}[t]
\centering
\caption{Cost-model sensitivity: geometric-mean $\Delta$ (\%) of
  \dSABRE$^{\dagger}$ (dS) and \BurstExt$^{\dagger}$ (dSE) vs.\
  \TeleSABRE{} as $c_{\mathrm{tele}}$ varies ($c_{\mathrm{swap}}=3$
  throughout).  Negative = better than \TeleSABRE{}.
  64q excludes Random (TS failure).}
\label{tab:sensitivity}
\setlength{\tabcolsep}{4pt}
\begin{tabular}{@{}r rr rr rr rr@{}}
\toprule
 & \multicolumn{2}{c}{25q} & \multicolumn{2}{c}{36q}
 & \multicolumn{2}{c}{25q$+$36q} & \multicolumn{2}{c}{64q} \\
\cmidrule(lr){2-3}\cmidrule(lr){4-5}\cmidrule(lr){6-7}\cmidrule(lr){8-9}
$c_{\mathrm{tele}}$ & dS & dSE & dS & dSE & dS & dSE & dS & dSE \\
\midrule
  3 & $+10.5$ & $\mathbf{-6.7}$  & $-12.3$ & $\mathbf{-17.6}$ & $-1.6$  & $\mathbf{-12.4}$ & $+18.3$ & $+9.8$  \\
  5 & $+9.3$  & $\mathbf{-8.0}$  & $-15.2$ & $\mathbf{-20.0}$ & $-3.8$  & $\mathbf{-14.2}$ & $+15.3$ & $+6.3$  \\
 10 & $+7.0$  & $\mathbf{-10.2}$ & $-20.6$ & $\mathbf{-24.4}$ & $-7.8$  & $\mathbf{-17.6}$ & $+9.8$  & $\mathbf{-0.1}$  \\
 20 & $+4.7$  & $\mathbf{-13.0}$ & $-26.9$ & $\mathbf{-29.6}$ & $-12.5$ & $\mathbf{-21.7}$ & $+3.3$  & $\mathbf{-7.7}$  \\
 30 & $+3.4$  & $\mathbf{-14.7}$ & $-30.5$ & $\mathbf{-32.5}$ & $-15.2$ & $\mathbf{-24.1}$ & $-0.5$  & $\mathbf{-12.1}$ \\
 50 & $+2.2$  & $\mathbf{-16.7}$ & $-34.5$ & $\mathbf{-35.7}$ & $-18.2$ & $\mathbf{-26.8}$ & $-4.6$  & $\mathbf{-17.0}$ \\
100 & $+1.1$  & $\mathbf{-18.8}$ & $-38.6$ & $\mathbf{-38.9}$ & $-21.2$ & $\mathbf{-29.6}$ & $-8.9$  & $\mathbf{-22.0}$ \\
\bottomrule
\end{tabular}
\end{table}

Three findings stand out.
First, \BurstExt$^{\dagger}$ strictly dominates \TeleSABRE{} in combined
cost on the 25-qubit and 36-qubit suites \emph{at every cost ratio tested},
including $c_{\mathrm{tele}}=c_{\mathrm{swap}}=3$ (i.e.\ treating a
teleportation as no more expensive than a local SWAP).
This means the EPR savings are large enough to overcome the slightly higher
SWAP count even in the most SWAP-friendly scenario.
Second, the 64-qubit crossover occurs between $c_{\mathrm{tele}}=5$ and
$c_{\mathrm{tele}}=10$: at the default $10{:}3$ ratio \BurstExt$^{\dagger}$
is essentially break-even ($-0.1\%$) with \TeleSABRE{} on this harder
suite, and strictly better for all larger ratios.
Since current photonic-link systems operate at ratios of ${\gg}30$~\cite{stephenson2020high},
\BurstExt{} is the right choice across the entire practically relevant range.
Third, \dSABRE$^{\dagger}$ without the wire-balanced extended set
exceeds \TeleSABRE{} in cost on the 25-qubit suite at all ratios
(though the gap narrows from $+10.5\%$ at $c_{\mathrm{tele}}=3$
to $+1.1\%$ at $c_{\mathrm{tele}}=100$): the dSABRE scoring heuristic
alone is insufficient to beat \TeleSABRE{}'s tighter SWAP discipline
on this suite, and the wire-balanced extended set is the decisive ingredient.

%%─────────────────────────────────────────────────────────────────────────────
\section{Related Work}
\label{sec:related}

\textbf{Distributed circuit compilation.}
Andres-Martinez and Heunen~\cite{andres2019automated} formalised the distributed
routing problem and proposed a hypergraph-partitioning approach.
Ferrari et al.~\cite{ferrari2021compiler} and Wu et al.~\cite{wu2022collcomm}
developed \SABRE{}-based heuristics.
Rodrigues et al.~\cite{rodrigues2022noise} add noise awareness.
Baker et al.~\cite{baker2020time} target time-optimal scheduling.
Russo et al.~\cite{russo2025telesabre} recently proposed \TeleSABRE{}, which
extends the scoring function to jointly consider teledata and telegate operations,
and uses a contracted communication-graph energy to evaluate inter-core candidates.
None of these approaches exploit commutativity in teleportation scoring.

\textbf{Commutativity in circuit compilation.}
Hong et al.~\cite{hong2024commutativity} introduced the EigenCluster Hypergraph
for intra-chip routing, an explicit data structure that partitions wires into
runs of mutually commuting gates.
The wire-balanced extended set developed in this paper takes a different
route: it captures commutativity opportunities implicitly through the
lookahead window, without ever computing eigen-directions or building a
hypergraph.
Itoko et al.~\cite{itoko2020optimization} and Sivarajah et
al.~\cite{sivarajah2020tket} exploit commutativity for gate cancellation and
depth reduction---orthogonal to EPR minimisation.
Vandersypen et al.~\cite{iten2016introduction} use eigen-directions in circuit
synthesis.

\textbf{SWAP minimisation.}
Cowtan et al.~\cite{cowtan2019phase}, Nash et al.~\cite{nash2020optimal}, and
Zulehner et al.~\cite{zulehner2018efficient} study SWAP minimisation on
monolithic chips.
Our work adapts the scoring philosophy to the distributed setting where the
dominant cost is EPR pairs rather than SWAPs.

%%─────────────────────────────────────────────────────────────────────────────
\section{Conclusion}
\label{sec:conclusion}

We presented \dSABRE{}, a distributed quantum circuit router with a five-term
teleportation scoring function and a proactive congestion-relief mechanism,
together with \BurstExt{}, a $\sim$50-line subclass that replaces the
topological-order extended set with a wire-balanced round-robin
construction.
Across twelve MQT-Bench circuits (six 25-qubit, six 36-qubit) on a
$2\!\times\!2$ distributed architecture, \dSABRE{} alone achieves a
geometric-mean 25.3\% EPR reduction over the C++ \TeleSABRE{} baseline
under identical initial layouts, and \BurstExt{} extends this to 33.6\%.
The wire-balanced extended set captures commutativity benefits without any
new scoring term, hyperparameter, or auxiliary data structure, and adds
essentially no compile-time overhead over \dSABRE{}.
We also showed that \dSABRE{} completes routing on a 64-qubit Random circuit
where \TeleSABRE{} fails to converge.

\textbf{Future directions.}
(i)~A C++ port of \dSABRE{} and \BurstExt{} to close the compile-time gap
with \TeleSABRE{}.
(ii)~Telegate (remote CX) operation support, which currently benefits
\TeleSABRE{} on QNN.
(iii)~Systematic evaluation on larger architectures (4--16 cores) and
deeper circuits.
(iv)~Integration with noise-aware cost models where EPR fidelity varies by link.
(v)~An explicit-commutativity routing variant built on a hypergraph
partitioning of each wire, addressed in a companion paper.

%%─────────────────────────────────────────────────────────────────────────────
\bibliographystyle{IEEEtran}
\bibliography{burst_sabre}

\end{document}
